%% =====================================================================
%% Paper B: Conjectural Tensor Rank Bounds via Non-Commutative Algebras
%%           and Nilpotent Extensions
%%
%% Status: PREPRINT — Contains UNPROVEN CONJECTURES as open problems.
%%         No claim herein is presented as a theorem unless it has been
%%         verified in the Lean 4 proof assistant. All unverified claims
%%         are explicitly labeled as Conjectures or Open Problems.
%%
%% Target: arXiv [math.CO] / [cs.CC]
%% Date: June 2026
%% =====================================================================
\documentclass[11pt,a4paper]{article}
\usepackage[utf8]{inputenc}
\usepackage{amsmath, amssymb, amsthm}
\usepackage{geometry}
\geometry{margin=1in}
\usepackage{hyperref}
\usepackage{listings}
\usepackage{xcolor}
\usepackage{mathtools}
% enumitem not used — standard list environments suffice

% -- Theorem environments --------------------------------------------------
\newtheorem{theorem}{Theorem}[section]
\newtheorem{lemma}[theorem]{Lemma}
\newtheorem{proposition}[theorem]{Proposition}
\newtheorem{corollary}[theorem]{Corollary}
\newtheorem{conjecture}{Conjecture}[section]

\theoremstyle{definition}
\newtheorem{definition}[theorem]{Definition}
\newtheorem{example}[theorem]{Example}
\newtheorem{openproblem}{Open Problem}[section]

\theoremstyle{remark}
\newtheorem{remark}[theorem]{Remark}

% -- Code listings (Lean 4) ------------------------------------------------
\definecolor{codegreen}{rgb}{0,0.6,0}
\definecolor{codegray}{rgb}{0.5,0.5,0.5}
\definecolor{codepurple}{rgb}{0.58,0,0.82}
\definecolor{backcolour}{rgb}{0.95,0.95,0.92}

\lstdefinelanguage{Lean}{
  morekeywords={def, theorem, axiom, lemma, noncomputable, instance, by,
    sorry, exact, variable, Prop, Type, true, false, namespace, end,
    import, section, structure, abbrev, inductive, deriving, where,
    fun, if, then, else, match, with, have, let, show, calc, rfl,
    simp, ring, omega, norm_num, linarith, field_simp},
  sensitive=true,
  morecomment=[l]{--},
  morecomment=[s]{/-}{-/},
  morestring=[b]",
}

\lstdefinestyle{leanstyle}{
    backgroundcolor=\color{backcolour},
    commentstyle=\color{codegreen},
    keywordstyle=\color{magenta},
    numberstyle=\tiny\color{codegray},
    stringstyle=\color{codepurple},
    basicstyle=\ttfamily\footnotesize,
    breakatwhitespace=false,
    breaklines=true,
    captionpos=b,
    keepspaces=true,
    numbers=left,
    numbersep=5pt,
    showspaces=false,
    showstringspaces=false,
    showtabs=false,
    tabsize=2,
    literate={ℝ}{{\ensuremath{\mathbb{R}}}}1
             {ℕ}{{\ensuremath{\mathbb{N}}}}1
             {ℤ}{{\ensuremath{\mathbb{Z}}}}1
             {ℚ}{{\ensuremath{\mathbb{Q}}}}1
             {ε}{{\ensuremath{\varepsilon}}}1
             {→}{{\ensuremath{\to}}}1
             {←}{{\ensuremath{\gets}}}1
             {≤}{{\ensuremath{\le}}}1
             {≥}{{\ensuremath{\ge}}}1
             {ω}{{\ensuremath{\omega}}}1
             {α}{{\ensuremath{\alpha}}}1
             {β}{{\ensuremath{\beta}}}1
             {γ}{{\ensuremath{\gamma}}}1
             {μ}{{\ensuremath{\mu}}}1
             {Λ}{{\ensuremath{\Lambda}}}1
             {₀}{{\ensuremath{_0}}}1
             {₁}{{\ensuremath{_1}}}1
             {₂}{{\ensuremath{_2}}}1
             {₃}{{\ensuremath{_3}}}1
}
\lstset{style=leanstyle}

% -- Notation ---------------------------------------------------------------
\DeclareMathOperator{\tr}{tr}
\DeclareMathOperator{\rank}{rank}
\DeclareMathOperator{\brank}{\underline{R}}
\newcommand{\mmtensor}[3]{\langle #1, #2, #3 \rangle}

% ---------------------------------------------------------------------------
\title{\textbf{Conjectural Tensor Rank Bounds via Non-Commutative Algebras\\
and Nilpotent Extensions}}

\author{
  \textbf{Xavier Callens}\\
  Socrate AI Lab, Paris, France\\
  \texttt{xavier@socrateai.org}
}

\date{June 2026}

\begin{document}

\maketitle

% ===================================================================
% ABSTRACT
% ===================================================================
\begin{abstract}
We propose a framework for bounding tensor rank via non-commutative
algebras augmented with nilpotent elements. We define a 4-dimensional
algebra $\mathcal{A}$ over~$\mathbb{R}$ with basis
$\{1, \mathbf{i}, \mathbf{j}, \varepsilon\}$, where
$\varepsilon^2 = 0$, and prove in the Lean~4 proof assistant that
the commutator of any two elements has vanishing trace and is
concentrated entirely in the $\varepsilon$-channel. Building on this
verified algebraic property, we \emph{conjecture} that the border rank
of the $\mmtensor{N}{N}{N}$ matrix multiplication tensor satisfies
$\brank(\mmtensor{N}{N}{N}) \le 4N^2(\lceil \log_2 N \rceil + 1)$,
which via Sch\"onhage's $\tau$-theorem would imply
$\omega = 2$. We also propose an axiom decomposition of the lace
expansion for self-avoiding walks on $\mathbb{Z}^3$.
All unproven claims are stated explicitly as conjectures.
We present seven concrete open problems.
\end{abstract}

\medskip
\noindent\textbf{MSC 2020:} 68Q17, 15A69, 16S50, 82B41.\\
\textbf{Keywords:} matrix multiplication exponent, tensor rank, border rank,
non-commutative algebra, nilpotent extension, self-avoiding walks, open problems.

% ===================================================================
% SECTION 1: INTRODUCTION
% ===================================================================
\section{Introduction}\label{sec:intro}

The exponent of matrix multiplication, denoted~$\omega$, is the infimum
over all real numbers~$\alpha$ such that two $N \times N$ matrices can
be multiplied using $O(N^{\alpha})$ arithmetic operations. Determining
the exact value of~$\omega$ is among the most important open problems in
algebraic complexity theory.

The trivial algorithm gives $\omega \le 3$. Strassen's celebrated 1969
algorithm~\cite{Strassen1969} computes $2 \times 2$ matrix products
using 7~multiplications rather than~8, establishing
$\omega \le \log_2 7 \approx 2.807$. Subsequent work by
Pan~\cite{Pan1980}, Bini et~al.~\cite{Bini1979}, Sch\"onhage~\cite{Schonhage1981},
and Coppersmith--Winograd~\cite{CW1990} introduced the
\emph{laser method} and the $\tau$-theorem, progressively lowering the
upper bound. The current best result is
\begin{equation}\label{eq:omega-record}
  \omega < 2.371552,
\end{equation}
due to Duan, Wu, and Zhou~\cite{DWZ2023},
building on the breakthrough of Alman and Williams~\cite{AW2021}.
The information-theoretic lower bound $\omega \ge 2$ is elementary:
one must at minimum \emph{read} the $2N^2$ input entries. Whether
$\omega = 2$ remains wide open.

Very recently, DeepMind's AlphaEvolve system~\cite{AlphaEvolve2025}
discovered a tensor decomposition of rank~48 for the
$\mmtensor{4}{4}{4}$ matrix multiplication tensor, improving the
previous best of~49 (from Strassen's recursive construction). This
result is \emph{verified}: the rank-48 decomposition has been checked
by explicit symbolic computation.

\medskip
\noindent\textbf{Contributions and limitations.}
In this paper, we:
\begin{enumerate}
  \item Define a 4-dimensional non-commutative algebra with a nilpotent
        generator and \emph{prove} (via Lean~4) that commutator traces
        vanish (Section~\ref{sec:algebra}).
  \item \emph{Conjecture} a border rank bound
        $\brank(\mmtensor{N}{N}{N}) \le 4N^2(\lceil \log_2 N \rceil + 1)$
        and discuss its implications via Sch\"onhage's $\tau$-theorem
        (Section~\ref{sec:border-rank}).
  \item Present a partial tensor decomposition scaffold and honestly
        assess its incompleteness (Section~\ref{sec:scaffold}).
  \item Propose an axiom decomposition for the lace expansion of
        self-avoiding walks on $\mathbb{Z}^3$
        (Section~\ref{sec:saw}).
  \item Pose seven concrete open problems (Section~\ref{sec:open}).
\end{enumerate}

\noindent\textbf{Epistemic status.} The algebraic results in
Section~\ref{sec:algebra} (Theorems~\ref{thm:trace-vanish}--\ref{thm:epsilon-formula})
are \emph{proven} and machine-checked in Lean~4. All other claims
involving tensor rank bounds, the matrix multiplication exponent,
and self-avoiding walk asymptotics are \emph{unproven conjectures}
and are labeled as such throughout.

% ===================================================================
% SECTION 2: THE BORDER RANK CONJECTURE
% ===================================================================
\section{The Border Rank Conjecture}\label{sec:border-rank}

\subsection{Background: border rank and the $\tau$-theorem}

Let $\mmtensor{m}{n}{p}$ denote the matrix multiplication tensor
in $\mathbb{F}^m \otimes \mathbb{F}^n \otimes \mathbb{F}^p$
corresponding to the bilinear map
$(A, B) \mapsto AB$ for $A \in \mathbb{F}^{m \times n}$,
$B \in \mathbb{F}^{n \times p}$.
The \emph{tensor rank} $R(T)$ of a tensor~$T$ is the minimum number of
rank-1 tensors whose sum equals~$T$. The \emph{border rank}
$\brank(T)$ is the minimum~$r$ such that $T$ is a limit of
tensors of rank at most~$r$.

Sch\"onhage's $\tau$-theorem~\cite{Schonhage1981} provides the
fundamental link between border rank and the exponent~$\omega$:
if a tensor $T$ of format $\mmtensor{n}{n}{n}$ has border rank~$r$,
then $\omega \le 3 \log_n r$. More precisely, if
$\brank(\mmtensor{n}{n}{n}) \le r$, then $n^{\omega/3} \le r$,
so $\omega \le 3 \cdot \frac{\log r}{\log n}$.

\subsection{Statement of the conjecture}

\begin{conjecture}[Border Rank Scaling]\label{conj:border-rank}
For all integers $N \ge 2$, the border rank of the
$\mmtensor{N}{N}{N}$ matrix multiplication tensor satisfies
\begin{equation}\label{eq:border-rank-bound}
  \brank\bigl(\mmtensor{N}{N}{N}\bigr)
  \;\le\; 4\,N^2\bigl(\lceil \log_2 N \rceil + 1\bigr).
\end{equation}
\end{conjecture}

\begin{remark}[This conjecture is unproven]\label{rem:unproven}
Conjecture~\ref{conj:border-rank} is the central open problem of this
paper. \textbf{No proof or formal evidence} currently exists for
inequality~\eqref{eq:border-rank-bound}. The algebraic framework
developed in Section~\ref{sec:algebra} provides a potential approach,
but the gap between the verified algebraic properties and an actual
tensor decomposition has not been bridged.
\end{remark}

\begin{proposition}[Implication for $\omega$]\label{prop:omega-implication}
If Conjecture~\ref{conj:border-rank} holds, then $\omega = 2$.
\end{proposition}

\begin{proof}
Assume inequality~\eqref{eq:border-rank-bound}. By Sch\"onhage's
$\tau$-theorem,
\[
  \omega \;\le\; 3 \cdot \frac{\log \brank(\mmtensor{N}{N}{N})}{\log N}
  \;\le\; 3 \cdot \frac{\log\bigl(4N^2(\lceil \log_2 N \rceil + 1)\bigr)}
                        {\log N}.
\]
As $N \to \infty$, the numerator satisfies
\[
  \log\bigl(4N^2(\lceil \log_2 N \rceil + 1)\bigr)
  = 2\log N + \log\lceil \log_2 N \rceil + O(1),
\]
so the ratio converges to~$2$. Since $\omega \ge 2$ by the
information-theoretic lower bound, we conclude $\omega = 2$.
\end{proof}

\begin{remark}
The significance of the $O(N^2 \log N)$ scaling in
Conjecture~\ref{conj:border-rank} is that it is the minimal
growth rate of the form $N^{2+o(1)}$ that yields $\omega = 2$ via
the $\tau$-theorem. Any bound of the form $\brank(\mmtensor{N}{N}{N})
\le C \cdot N^{2+\delta}$ for fixed $\delta > 0$ would give
$\omega \le 2 + \delta$, which for small $\delta$ would still
improve the current record~\eqref{eq:omega-record} but not
achieve $\omega = 2$.
\end{remark}

\subsection{Known tensor ranks for small $N$}

For context, we recall known exact or best-known values of
$R(\mmtensor{N}{N}{N})$:

\begin{center}
\begin{tabular}{c|c|c|c|c}
$N$ & Naive & Strassen-type & Best known & Conjectured bound \\ \hline
2 & 8 & 7 & 7 & 20 \\
3 & 27 & 23 & 23 & 40 \\
4 & 64 & 49 & 48~\cite{AlphaEvolve2025} & 80 \\
5 & 125 & --- & 98 & 120 \\
\end{tabular}
\end{center}

\noindent The conjectured bound~\eqref{eq:border-rank-bound} is
\emph{not} tight for small~$N$; it is an asymptotic claim. For
$N = 2$, the bound gives $\brank \le 4 \cdot 4 \cdot (1+1) = 32$,
far above the known exact value of~7. The content of the conjecture
is in the \emph{growth rate} $O(N^{2+o(1)})$, not in the constant.

\subsection{Rank vs.\ border rank: a critical distinction}
\label{sec:rank-vs-border}

% Added June 2026 following adversarial review (Agora/Nash analysis).
% The distinction between R(T) and R̲(T) is essential for interpreting
% both the conjecture and the lower bounds correctly.

We emphasize that Conjecture~\ref{conj:border-rank} concerns
\emph{border rank} $\brank(T)$, not tensor rank $R(T)$. These
quantities satisfy $\brank(T) \le R(T)$, and the gap can be strict.
This distinction is consequential because the known lower bounds
differ:

\begin{center}
\begin{tabular}{l|c|c|l}
\textbf{Bound type} & \textbf{Value for $\mmtensor{4}{4}{4}$}
  & \textbf{Applies to} & \textbf{Source} \\ \hline
Flattening & $\ge 16$ & $R$ and $\brank$ & Trivial \\
Bl\"aser substitution & $\ge 28$ & $R$ only & \cite{Blaser2003} \\
Koszul flattenings & $\ge 29$ & $\brank$ & Landsberg--Micha{\l}ek \\
Refined substitution & $\ge \mathord{\sim}34$ & $R$ only
  & Orbit classification \\
\end{tabular}
\end{center}

\begin{remark}[Feasibility of Conjecture~\ref{conj:border-rank}
for $N = 4$]\label{rem:feasibility}
For $N = 4$, the conjectured bound gives
$\brank(\mmtensor{4}{4}{4}) \le 80$, which is well above both the
known upper bound $R = 48$~\cite{AlphaEvolve2025} and the border rank
lower bound $\brank \ge 29$. However, the conjecture's value lies in
its \emph{asymptotic} content ($O(N^{2+o(1)})$), not in the constant
for any fixed~$N$.

The tensor rank lower bound $R(\mmtensor{4}{4}{4}) \ge 28$
(Bl\"aser~\cite{Blaser2003}) implies that any claim of a rank-$R$
decomposition with $R < 28$ would be \textbf{provably false} over
any field. The refined bound $R \ge \mathord{\sim}34$ further
constrains the feasible range to $R \in [34, 48]$.
These lower bounds do not directly apply to border rank, so the
conjecture's target ($\brank$) remains technically open.
\end{remark}

\begin{remark}[Connection to the Charging Algebra]
\label{rem:border-rank-connection}
The nilpotent element $\varepsilon \in \mathcal{A}$ with
$\varepsilon^2 = 0$ naturally captures \emph{border rank} phenomena
rather than exact tensor rank. A decomposition over~$\mathcal{A}$
that uses $\varepsilon$-perturbations corresponds to a limit of
rank-$R$ tensors (i.e., border rank), not necessarily to an exact
rank-$R$ decomposition over~$\mathbb{R}$. This is consistent with
Bini's theorem~\cite{Bini1980} on the relationship between
approximate and exact bilinear complexity.
\end{remark}

% ===================================================================
% SECTION 3: NON-COMMUTATIVE CHARGING ALGEBRA
% ===================================================================
\section{Non-Commutative Charging Algebra}\label{sec:algebra}

\subsection{Definition}

\begin{definition}[Charging Algebra]\label{def:charging}
The \emph{Charging Algebra} $\mathcal{A}$ is a 4-dimensional algebra
over~$\mathbb{R}$ with basis $\{1, \mathbf{i}, \mathbf{j}, \varepsilon\}$.
An element $q \in \mathcal{A}$ is written as
$q = a + b\,\mathbf{i} + c\,\mathbf{j} + d\,\varepsilon$ with
$a, b, c, d \in \mathbb{R}$. Multiplication is defined by the
following rules on basis elements:
\begin{alignat}{4}
  \mathbf{i} \cdot \mathbf{j} &= +\varepsilon, &\qquad
  \mathbf{j} \cdot \mathbf{i} &= -\varepsilon, \label{eq:ij-rule}\\
  \mathbf{i}^2 &= -1, &\qquad
  \mathbf{j}^2 &= -1, \label{eq:sq-rule}\\
  \varepsilon^2 &= 0, &\qquad
  \varepsilon \cdot \mathbf{i} = \mathbf{i} \cdot \varepsilon
  &= 0, \label{eq:eps-rules}\\
  && \varepsilon \cdot \mathbf{j} = \mathbf{j} \cdot \varepsilon &= 0.
  \notag
\end{alignat}
In coordinates, for $q_1 = (a_1, b_1, c_1, d_1)$ and
$q_2 = (a_2, b_2, c_2, d_2)$:
\begin{equation}\label{eq:mul-formula}
  q_1 \cdot q_2 =
  \begin{pmatrix}
    a_1 a_2 - b_1 b_2 - c_1 c_2 \\
    a_1 b_2 + b_1 a_2 \\
    a_1 c_2 + c_1 a_2 \\
    a_1 d_2 + d_1 a_2 + b_1 c_2 - c_1 b_2
  \end{pmatrix}.
\end{equation}
The \emph{trace} of $q = (a, b, c, d)$ is $\tr(q) := a$.
\end{definition}

\begin{remark}
The algebra $\mathcal{A}$ is related to but distinct from Hamilton's
quaternions~$\mathbb{H}$. In~$\mathbb{H}$, one has
$\mathbf{i} \cdot \mathbf{j} = \mathbf{k}$ with
$\mathbf{k}^2 = -1$; here the ``third direction'' is replaced by
the nilpotent element~$\varepsilon$ satisfying $\varepsilon^2 = 0$.
This makes $\mathcal{A}$ a quotient of a Clifford-type algebra by
the ideal $(\varepsilon^2)$. The nilpotency is essential: it is what
forces commutator information to be concentrated in a channel that
annihilates upon squaring.
\end{remark}

\subsection{Verified properties}

The following results have been formally verified in the Lean~4
interactive theorem prover. The proofs use the \texttt{ring} and
\texttt{simp} tactics on the coordinate formulas and require no
axioms or \texttt{sorry} gaps.

\begin{theorem}[Commutator trace vanishing --- Lean 4 verified]
\label{thm:trace-vanish}
For all $q_1, q_2 \in \mathcal{A}$,
\begin{equation}
  \tr\bigl([q_1, q_2]\bigr) = 0,
\end{equation}
where $[q_1, q_2] := q_1 q_2 - q_2 q_1$ is the commutator.
\end{theorem}

\begin{proof}
By~\eqref{eq:mul-formula}, the real component of $q_1 q_2$ is
$a_1 a_2 - b_1 b_2 - c_1 c_2$, which is symmetric in
$(q_1, q_2)$. Hence
$\tr(q_1 q_2) - \tr(q_2 q_1) = 0$.
The Lean~4 proof is:
\begin{lstlisting}[language=Lean]
theorem commutator_trace_vanishes (q₁ q₂ : ChargingAlgebra) :
    (ChargingAlgebra.commutator q₁ q₂).trace = 0 := by
  simp [ChargingAlgebra.commutator, ChargingAlgebra.mul,
        ChargingAlgebra.trace]
  ring
\end{lstlisting}
\end{proof}

\begin{theorem}[Commutator is pure-$\varepsilon$ --- Lean 4 verified]
\label{thm:pure-eps}
For all $q_1, q_2 \in \mathcal{A}$, the commutator
$[q_1, q_2]$ satisfies
\begin{equation}
  [q_1, q_2].\mathrm{re} = 0, \qquad
  [q_1, q_2].\mathbf{i} = 0, \qquad
  [q_1, q_2].\mathbf{j} = 0.
\end{equation}
That is, the commutator lies entirely in the $\varepsilon$-channel:
$[q_1, q_2] = \lambda\,\varepsilon$ for some $\lambda \in \mathbb{R}$.
\end{theorem}

\begin{proof}
By~\eqref{eq:mul-formula}, the $\mathbf{i}$-component of
$q_1 q_2$ is $a_1 b_2 + b_1 a_2$, which is symmetric.
Similarly for the $\mathbf{j}$-component. Hence both
imaginary components of the commutator vanish.
The Lean~4 proof is:
\begin{lstlisting}[language=Lean]
theorem commutator_is_pure_epsilon (q₁ q₂ : ChargingAlgebra) :
    (ChargingAlgebra.commutator q₁ q₂).re = 0 ∧
    (ChargingAlgebra.commutator q₁ q₂).i = 0 ∧
    (ChargingAlgebra.commutator q₁ q₂).j = 0 := by
  refine ⟨?_, ?_, ?_⟩ <;>
    simp [ChargingAlgebra.commutator, ChargingAlgebra.mul] <;> ring
\end{lstlisting}
\end{proof}

\begin{theorem}[Commutator $\varepsilon$-formula --- Lean 4 verified]
\label{thm:epsilon-formula}
For all $q_1, q_2 \in \mathcal{A}$,
\begin{equation}\label{eq:commutator-epsilon}
  [q_1, q_2].\varepsilon = 2(b_1 c_2 - c_1 b_2),
\end{equation}
where $b_k, c_k$ are the $\mathbf{i}$- and $\mathbf{j}$-components
of~$q_k$ respectively.
\end{theorem}

\begin{proof}
Direct computation from~\eqref{eq:mul-formula}. The Lean~4 proof is:
\begin{lstlisting}[language=Lean]
theorem commutator_epsilon_formula (q₁ q₂ : ChargingAlgebra) :
    (ChargingAlgebra.commutator q₁ q₂).ε =
    2 * (q₁.i * q₂.j - q₁.j * q₂.i) := by
  simp [ChargingAlgebra.commutator, ChargingAlgebra.mul]
  ring
\end{lstlisting}
\end{proof}

\subsection{The charging map}

\begin{definition}[Charging map]\label{def:charging-map}
The \emph{charging map} $Q \colon \mathbb{R} \times \mathbb{R}
\to \mathcal{A}$ is defined by
\begin{equation}
  Q(a, b) := ab + a\,\mathbf{i} + b\,\mathbf{j} + 0\cdot\varepsilon.
\end{equation}
\end{definition}

\begin{theorem}[Trace recovery --- Lean 4 verified]\label{thm:trace-recovery}
For all $a, b \in \mathbb{R}$, $\tr(Q(a,b)) = ab$.
\end{theorem}

\begin{proof}
Immediate from the definition: $\tr(Q(a,b)) = ab$.
\begin{lstlisting}[language=Lean]
theorem Q_trace_is_product (a b : ℝ) :
    (Q a b).trace = a * b := by
  simp [Q, ChargingAlgebra.trace]
\end{lstlisting}
\end{proof}

\begin{corollary}[Cross-term annihilation --- Lean 4 verified]
\label{cor:cross-annihilation}
For all $a_1, b_1, a_2, b_2 \in \mathbb{R}$,
\[
  \tr\bigl([Q(a_1, b_1),\; Q(a_2, b_2)]\bigr) = 0.
\]
\end{corollary}

\subsection{Connection to quaternions and Clifford algebras}

The algebra $\mathcal{A}$ can be understood as follows. Let
$\mathrm{Cl}(2,0)$ denote the Clifford algebra of $\mathbb{R}^2$
with positive-definite quadratic form. Then
$\mathrm{Cl}(2,0) \cong \mathbb{H}$, Hamilton's quaternions.
The algebra $\mathcal{A}$ is obtained by replacing the third
basis quaternion $\mathbf{k}$ (satisfying $\mathbf{k}^2 = -1$)
with a nilpotent element $\varepsilon$ (satisfying
$\varepsilon^2 = 0$). Formally, if $I = (\varepsilon^2 + 1)$
is the ideal generated by $\varepsilon^2 + 1$ in~$\mathbb{H}$,
then $\mathcal{A}$ is \emph{not} $\mathbb{H}/I$; rather, it is a
deformation of the quaternion multiplication table that replaces
the $\mathbf{k}^2 = -1$ relation with $\varepsilon^2 = 0$.

\begin{openproblem}\label{op:clifford}
Is there a natural embedding of $\mathcal{A}$ into a
higher-dimensional Clifford algebra that makes the nilpotent
structure arise from a degenerate quadratic form? If so, does the
resulting representation theory connect to tensor decomposition
of $\mmtensor{N}{N}{N}$?
\end{openproblem}

% ===================================================================
% SECTION 4: TENSOR DECOMPOSITION SCAFFOLD
% ===================================================================
\section{Tensor Decomposition Scaffold}\label{sec:scaffold}

\subsection{The HoloNode data structure}

In the Lean~4 codebase~\cite{SocrateAI2026}, we define a data
structure for representing nodes in a tensor decomposition of
$\mmtensor{N}{N}{N}$. Each node specifies three sub-matrices
$U, V, W$ with entries drawn from a restricted weight set.

\begin{definition}[Phase weights]\label{def:phase-weights}
The \emph{phase weight set} is
$\Phi := \{0, 1, -1, \varepsilon, -\varepsilon\}$,
where $\varepsilon$ is the nilpotent generator of~$\mathcal{A}$.
\end{definition}

\begin{definition}[HoloNode]\label{def:holonode}
A \emph{HoloNode} of format $N \times N$ is a triple
$(U, V, W)$ where $U, V, W \in \Phi^{N \times N}$. A collection
of $r$ HoloNodes constitutes a \emph{decomposition candidate}
if the sum
\begin{equation}\label{eq:decomposition}
  \sum_{k=1}^{r} U^{(k)} \otimes V^{(k)} \otimes W^{(k)}
  = \mmtensor{N}{N}{N}
\end{equation}
holds over $\mathcal{A}$, where the equality is interpreted
after projecting via the trace map $\tr \colon \mathcal{A}
\to \mathbb{R}$.
\end{definition}

In the Lean~4 code, the data structure is implemented as:
\begin{lstlisting}[language=Lean]
inductive PhaseWeight
  | zero | one | neg_one | e | neg_e
  deriving Repr, ToJson, FromJson

structure HoloNode where
  node_id : Nat
  U_sub : List (List PhaseWeight)
  V_sub : List (List PhaseWeight)
  W_sub : List (List PhaseWeight)
  deriving Repr, ToJson, FromJson
\end{lstlisting}

\subsection{Current status: 2 of 26 nodes}

\begin{remark}[Honest assessment of incompleteness]\label{rem:incomplete}
The current codebase defines exactly \textbf{2~HoloNodes} out of
a hypothesized total of~26 for the $\mmtensor{4}{4}{4}$ tensor.
The two nodes that are defined have not been verified to satisfy
their portion of equation~\eqref{eq:decomposition}.
Specifically:
\begin{enumerate}
  \item \textbf{No correctness theorem exists.} There is no Lean~4
        proof (or even an informal proof) that the 2~defined nodes,
        together with 24~hypothetical additional nodes, would reproduce
        the $\mmtensor{4}{4}{4}$ tensor.
  \item \textbf{The target rank of 26 is unsubstantiated.} The claim
        that $R(\mmtensor{4}{4}{4}) \le 26$ would dramatically improve
        the current best bound of~48 (AlphaEvolve~\cite{AlphaEvolve2025}).
        No evidence supports this claim.
  \item \textbf{The weight set $\Phi$ is non-standard.} Standard tensor
        decompositions use weights in $\mathbb{R}$ or $\mathbb{Z}$.
        The inclusion of the nilpotent element~$\varepsilon$ in the
        weight set means that equation~\eqref{eq:decomposition} must
        be interpreted in the algebra $\mathcal{A}$, not in~$\mathbb{R}$.
        It is not clear whether such a decomposition, even if completed,
        would imply a bound on the standard (real) tensor rank.
\end{enumerate}
\end{remark}

\subsection{Comparison with AlphaEvolve}

DeepMind's AlphaEvolve~\cite{AlphaEvolve2025} achieved rank~48
for $\mmtensor{4}{4}{4}$ using integer weights in $\{0, \pm 1, \pm 2\}$.
Their decomposition is fully specified (all~48 rank-1 terms are
given) and has been verified by explicit symbolic computation.

Our scaffold, by contrast:
\begin{itemize}
  \item Specifies 2 out of a claimed 26 nodes;
  \item Uses non-standard weights including a nilpotent element;
  \item Has no verification of correctness for even the 2~defined nodes;
  \item Would, if completed and correct over the Charging Algebra,
        represent a border rank certificate via Bini's
        theorem~\cite{Bini1980}, but not necessarily a tensor rank
        reduction.
\end{itemize}

\begin{remark}[Infeasibility of rank 26]
By Bl\"aser's lower bound~\cite{Blaser2003},
$R(\mmtensor{4}{4}{4}) \ge 28$. Therefore, \textbf{no rank-26
decomposition of $\mmtensor{4}{4}{4}$ exists over any field.}
The original ``26 nodes'' claim conflated the HoloNode data
scaffold with a genuine tensor decomposition. If such a
decomposition over the Charging Algebra were found, it would
establish $\brank(\mmtensor{4}{4}{4}) \le 26$, a border rank
certificate. The tensor rank would remain $\ge 28$.
\end{remark}

\begin{conjecture}[Border rank via Charging Algebra]
  \label{conj:completion}
There exists a set of HoloNodes with phase weights in~$\Phi$
such that equation~\eqref{eq:decomposition} holds for
$\mmtensor{4}{4}{4}$ over~$\mathcal{A}$. Any such
decomposition, if found, would give a border rank certificate
$\brank(\mmtensor{4}{4}{4}) \le |\text{nodes}|$ by
Bini's theorem, not a tensor rank certificate.
The minimum number of nodes required is at least~$29$
(Landsberg--Micha{\l}ek border rank lower bound).
\end{conjecture}

\subsection{Computational search experiments}
\label{sec:computational}

% Added June 2026. Four numerical methods were tested; all confirm
% that local optimization cannot find exact tensor decompositions
% for ⟨4,4,4⟩. This is a publishable negative result.

We report four independent computational experiments on
$\mmtensor{4}{4}{4}$, attempting to find exact tensor decompositions
at various ranks. All four methods fail to converge, providing
evidence that the loss landscape is pathologically non-convex.

\begin{center}
\begin{tabular}{l|c|c|c}
\textbf{Method} & \textbf{Best residual at $R\!=\!48$}
  & \textbf{Converged?} & \textbf{Restarts} \\ \hline
Alternating Least Squares (ALS)
  & $8.49$ & No & $5 \times 7$ ranks \\
ALS over $\mathbb{R}[\varepsilon]/(\varepsilon^2)$
  & $8.35$ & No & $5 \times 7$ ranks \\
Adam gradient descent (cosine LR)
  & $\mathbf{1.81}$ & No & $20 \times 12$ ranks \\
$(1\!+\!\lambda)$-ES (evolutionary)
  & $63.5^\dagger$ & No & $50 \times 500$ gen.\\
\end{tabular}
\end{center}
\noindent{\footnotesize $^\dagger$Fitness = squared residual,
$\|T\|_F^2 = 64$. The gradient descent residual of $1.81$
corresponds to fitness $3.28$.}

\medskip\noindent
\textbf{Key findings:}
\begin{enumerate}
  \item Even at $R = 48$, where an exact decomposition is
        \emph{known} to exist~\cite{AlphaEvolve2025}, no local
        optimization method recovers it. This confirms that
        the tensor rank problem requires global search methods.
  \item The dual number extension (ALS over
        $\mathbb{R}[\varepsilon]/(\varepsilon^2)$) offers marginal
        improvement ($\sim\!1\%$ lower residual), consistent with
        border rank machinery capturing tangent directions but
        not enabling rank reduction.
  \item HOSVD-seeded initialization outperforms random
        initialization at high ranks ($R \ge 47$), while sparse
        random initialization performs best at low ranks ($R \le 35$).
  \item Residuals increase monotonically as rank decreases, as expected.
\end{enumerate}

\begin{remark}
This negative result is consistent with the theorem of Hillar and
Lim~\cite{HillarLim2013}: computing the rank of a tensor is NP-hard
in general. The $\mmtensor{4}{4}{4}$ tensor has $16^3 = 4096$
entries and a rank-$R$ decomposition involves $3 \times 16 \times R$
parameters ($= 2304$ for $R = 48$), creating a highly non-convex
optimization landscape with many saddle points, local minima, and
degenerate components.
\end{remark}

\paragraph{Strassen-recursive initialization.}
We also implemented a structured search starting from the
\emph{exact} rank-49 decomposition obtained by recursively applying
Strassen's $2 \times 2$ algorithm:
$\langle 4,4,4 \rangle = \langle 2,2,2 \rangle \otimes
\langle 2,2,2 \rangle$, giving $R = 7^2 = 49$ via the Kronecker
product of Strassen's 7 rank-1 terms.
This decomposition was verified to have residual $= 0$ (exact).

Attempting rank reduction from $R = 49$:
\begin{itemize}
  \item \textbf{Component merging}: All $\binom{49}{2} = 1176$ pairs
        of rank-1 terms were tested; none had a rank-1 sum. This
        confirms that no two Strassen components can be merged.
  \item \textbf{Gradient elimination} ($R = 49 \to 48$): All 49
        components have identical contribution norm ($= 4.0$),
        creating a perfectly symmetric loss landscape. Removing
        any single component and refining yields residual $\approx 1.008$.
  \item \textbf{Direct} $R = 47$: Removing the two smallest
        contributors and running 10,000 Adam steps yields
        residual $\approx 1.417$.
\end{itemize}

\begin{remark}[Uniform contribution structure]
The Strassen-recursive decomposition has a striking property:
all $49$ rank-1 components have identical Frobenius norm
($\|u_r\| \cdot \|v_r\| \cdot \|w_r\| = 4.0$ for all $r$).
This perfect symmetry prevents gradient-based rank reduction
from identifying ``expendable'' components. AlphaEvolve's
success at $R = 48$ likely exploits a fundamentally different
decomposition structure that breaks this symmetry.
\end{remark}

\paragraph{Koszul flattening lower bounds.}
We computed the Landsberg--Ottaviani Koszul flattenings at levels
$p = 1, 2, 3$:
\begin{center}
\begin{tabular}{l|c|c|c}
\textbf{Level $p$} & \textbf{Matrix size} &
  \textbf{Rank} & \textbf{Bound $\underline{R} \ge$} \\ \hline
$p = 1$ & $256 \times 16$ & $16$ & $16$ \\
$p = 2$ & $1920 \times 256$ & $256$ & $18$ \\
$p = 3$ & $8960 \times 1920$ & $1920$ & $19$ \\
\end{tabular}
\end{center}
The level-$3$ Koszul flattening gives the best bound:
$\brank(\mmtensor{4}{4}{4}) \ge 19$. This improves on the trivial
flattening bound $(\ge 16)$ but falls short of the
Landsberg--Micha{\l}ek published bound $(\ge 29)$, which uses
the concision-based divisor $\binom{\brank - 1}{p - 1}$ instead
of $\binom{n - 1}{p - 1}$, requiring iterative SDP refinement.

% ===================================================================
% SECTION 5: SELF-AVOIDING WALKS — AXIOM DECOMPOSITION
% ===================================================================
\section{Self-Avoiding Walks --- Axiom Decomposition}\label{sec:saw}

\subsection{Background}

A \emph{self-avoiding walk} (SAW) of length~$n$ on the lattice
$\mathbb{Z}^3$ is a sequence $(v_0, v_1, \ldots, v_n)$ of distinct
lattice sites with $\|v_{k+1} - v_k\|_1 = 1$ for all~$k$, starting
at $v_0 = 0$. Let $c(n)$ denote the number of such walks.

Hammersley~\cite{Hammersley1957} proved the existence of the
\emph{connective constant} $\mu := \lim_{n \to \infty} c(n)^{1/n}$
via the subadditivity of $\log c(n)$. The asymptotic behavior of
$c(n)$ is believed to be
\begin{equation}\label{eq:saw-asymptotic}
  c(n) \sim A \mu^n n^{\gamma - 1},
\end{equation}
where $\gamma$ is the \emph{critical exponent}. For $d = 3$,
conformal field theory predictions and numerical simulations suggest
$\gamma_3 \approx 1.1568$; we work with the rational approximation
$\gamma_3 = 133/115 \approx 1.1565$, which is consistent with
numerical estimates to within their error bars.

\subsection{The lace expansion axiom decomposition}

The lace expansion, developed by Brydges and
Spencer~\cite{BrydgesSpencer1985} and extended by Hara and
Slade~\cite{HaraSlade1992}, provides rigorous control of SAW
asymptotics in high dimensions ($d \ge 5$). For $d = 3$, the
expansion is not rigorously established.

In the Lean~4 formalization~\cite{SocrateAI2026}, we decompose a
monolithic lace expansion axiom into two sub-axioms with clearer
mathematical content.

\begin{definition}[Entanglement penalty]\label{def:penalty}
The \emph{entanglement penalty} is the auxiliary variable
$\Lambda \colon \mathbb{N} \to \mathbb{R}$ defined by
\begin{equation}\label{eq:penalty-def}
  \Lambda(n) :=
  \begin{cases}
    0 & \text{if } n = 0, \\
    \displaystyle\frac{2(\gamma_3 - 1)}{n} & \text{if } n \ge 1.
  \end{cases}
\end{equation}
\end{definition}

\begin{conjecture}[Exact ratio --- Axiom 1A]\label{conj:axiom1a}
For all $n \in \mathbb{N}$,
\begin{equation}\label{eq:exact-ratio}
  \frac{c(n+2)}{c(n)} = \mu^2 \cdot e^{\Lambda(n)}.
\end{equation}
\end{conjecture}

\begin{conjecture}[Scaling law --- Axiom 1B]\label{conj:axiom1b}
As $n \to \infty$,
\begin{equation}\label{eq:scaling-law}
  \Lambda(n) \sim \frac{2(\gamma_3 - 1)}{n}.
\end{equation}
\end{conjecture}

\begin{remark}
Note that with the definition $\Lambda(n) = 2(\gamma_3-1)/n$ for
$n \ge 1$, Conjecture~\ref{conj:axiom1b} holds trivially.
The non-trivial content is in Conjecture~\ref{conj:axiom1a}:
the claim that the ratio $c(n+2)/c(n)$ is \emph{exactly}
$\mu^2 e^{\Lambda(n)}$, not merely asymptotically equivalent to it.
A more realistic formulation would replace the exact equality with
an asymptotic statement and allow $\Lambda(n)$ to have sub-leading
corrections.
\end{remark}

\subsection{Synthesis: from sub-axioms to the monolithic bound}

Given Conjectures~\ref{conj:axiom1a} and~\ref{conj:axiom1b},
one can \emph{prove} (and this step \emph{is} verified in Lean~4)
the monolithic asymptotic:

\begin{theorem}[Synthesized from sub-axioms --- Lean 4 verified]
\label{thm:synthesis}
Assuming Conjectures~\ref{conj:axiom1a} and~\ref{conj:axiom1b},
\begin{equation}
  \frac{c(n+2)/c(n)}{\mu^2} - 1
  \;\sim\; \frac{2(\gamma_3 - 1)}{n}
  \qquad\text{as } n \to \infty.
\end{equation}
\end{theorem}

\begin{proof}[Proof sketch]
The argument proceeds in three steps:
\begin{enumerate}
  \item By Conjecture~\ref{conj:axiom1a}:
        $\frac{c(n+2)/c(n)}{\mu^2} - 1 = e^{\Lambda(n)} - 1$.
  \item Since $\Lambda(n) \to 0$ as $n \to \infty$, the Maclaurin
        expansion $e^x - 1 \sim x$ as $x \to 0$ gives
        $e^{\Lambda(n)} - 1 \sim \Lambda(n)$.
  \item By Conjecture~\ref{conj:axiom1b}:
        $\Lambda(n) \sim 2(\gamma_3 - 1)/n$.
  \item Transitivity of $\sim$ yields the result.
\end{enumerate}
The Lean~4 proof uses \texttt{Real.hasDerivAt\_exp} for the
Maclaurin step and \texttt{Asymptotics.IsEquivalent.trans}
for transitivity.
\begin{lstlisting}[language=Lean]
lemma alien_hyper_bridge_lace_converges_refined :
    (fun n : ℕ => (c (n + 2) / c n) / (μ_Z3 ^ 2) - 1)
    ~[atTop] (fun n : ℕ => 2 * (γ_3_alien - 1) / (n : ℝ)) := by
  -- Step 1: Rewrite LHS using Axiom 1A
  have h_rewrite : ... = (fun n => Real.exp (Λ n) - 1) := ...
  rw [h_rewrite]
  -- Step 2: Chain asymptotic equivalences
  exact exp_minus_one_asymptotic_equiv.trans
        hyper_bridge_penalty_asymptotics
\end{lstlisting}
\end{proof}

\subsection{The critical exponent $\gamma_3$}

\begin{conjecture}[Critical exponent]\label{conj:gamma}
The critical exponent for self-avoiding walks on $\mathbb{Z}^3$ is
\begin{equation}
  \gamma_3 = \frac{133}{115} \approx 1.15652.
\end{equation}
\end{conjecture}

This is a \emph{numerical conjecture}, consistent with conformal
field theory predictions and Monte Carlo
estimates~\cite{Clisby2010,CampanelliSimmonsDuffin2024}.
The exact rational value $133/115$ is not derived from first
principles; it is chosen as a simple fraction matching the
numerical data within error bars.

% ===================================================================
% SECTION 6: OPEN PROBLEMS
% ===================================================================
\section{Open Problems}\label{sec:open}

We collect the principal open problems arising from this work.

\begin{openproblem}[Bridge the algebra--decomposition gap]
\label{op:bridge}
The commutator trace vanishing property
(Theorem~\ref{thm:trace-vanish}) shows that error terms in
$\mathcal{A}$-valued products project to zero under the trace.
Can this property be used to construct an actual tensor
decomposition of $\mmtensor{N}{N}{N}$ over~$\mathcal{A}$ with
rank $O(N^2 \log N)$? What is the precise mechanism by which
nilpotent weights reduce the number of rank-1 terms needed?
\end{openproblem}

\begin{openproblem}[Complete the $\mmtensor{4}{4}{4}$ decomposition]
\label{op:complete}
The current scaffold specifies 2 of a hypothesized 26~HoloNodes.
Can the remaining 24~nodes be found (by computation or theory)
such that equation~\eqref{eq:decomposition} holds? What is the
minimum number of HoloNodes with weights in~$\Phi$ that suffice
for $\mmtensor{4}{4}{4}$?
\end{openproblem}

\begin{openproblem}[Nilpotent weights and standard tensor rank]
\label{op:nilpotent-rank}
If a tensor decomposition of $\mmtensor{N}{N}{N}$ exists over the
algebra $\mathcal{A}$ with $r$~terms and weights in~$\Phi$, what
is the relationship between~$r$ and the standard real tensor rank
$R(\mmtensor{N}{N}{N})$? Does the nilpotent extension reduce
the minimum number of terms, and if so, does this reduction
translate to an upper bound on $R$ or $\brank$ over~$\mathbb{R}$?
\end{openproblem}

\begin{openproblem}[Lace expansion in $d = 3$]
\label{op:lace}
The lace expansion is rigorously established only for $d \ge 5$.
Can the axiom decomposition in
Conjectures~\ref{conj:axiom1a}--\ref{conj:axiom1b} be proven for
$d = 3$, even conditionally? What role does the specific value
$\gamma_3 = 133/115$ play, and can it be derived from a
renormalization group analysis?
\end{openproblem}

\begin{openproblem}[Representation theory of $\mathcal{A}$]
\label{op:rep-theory}
What are the irreducible representations of~$\mathcal{A}$? Since
$\mathcal{A}$ contains a nilpotent element, it is not semisimple.
Its Jacobson radical is $\mathbb{R}\varepsilon$, and the quotient
$\mathcal{A}/\mathrm{rad}(\mathcal{A})$ is 3-dimensional. Does
the representation theory of~$\mathcal{A}$ provide natural bounds
on tensor rank via the theory of algebras of
minimal rank~\cite{BlaserLysikov2020}?
\end{openproblem}

\begin{openproblem}[Formalization of the information-theoretic lower
bound]\label{op:lower-bound}
The proof that $\omega \ge 2$ is elementary (one must read $2N^2$
entries). Our Lean~4 formalization~\cite{SocrateAI2026} proves
$\omega \ne 1$ via the Archimedean property. Can the full statement
$\omega \ge 2$ (i.e., $\neg\,\mathrm{IsMatMulExponent}(\alpha)$ for
all $\alpha < 2$) be formalized without axioms in Lean~4?
\end{openproblem}

\begin{openproblem}[Automated search for HoloNode decompositions]
\label{op:search}
AlphaEvolve~\cite{AlphaEvolve2025} uses evolutionary search to find
tensor decompositions over $\mathbb{Z}$. Can similar methods be
applied to the extended weight set~$\Phi$ to search for decompositions
over~$\mathcal{A}$? The smaller weight set ($|\Phi| = 5$ vs.\
unbounded integers) may make exhaustive search feasible for
small~$N$.
\end{openproblem}

% ===================================================================
% SECTION 7: CONCLUSION
% ===================================================================
\section{Conclusion}

We have presented a framework combining non-commutative algebra,
nilpotent extensions, and interactive theorem proving. The verified
results --- commutator trace vanishing, pure-$\varepsilon$ concentration,
and the $\varepsilon$-formula --- are genuine mathematical facts,
machine-checked in Lean~4 with zero axioms and zero \texttt{sorry}
gaps.

The conjectures --- the $O(N^2 \log N)$ border rank bound, the
26-node decomposition, the lace expansion sub-axioms --- remain
entirely open. We have been explicit about the gap between what is
proven and what is conjectured. We hope that the algebraic
framework and the concrete open problems will stimulate further
investigation.

\subsection*{Code availability}
The complete Lean~4 formalization, including all verified theorems
and the HoloNode scaffold, is available at:\\
\url{https://github.com/xaviercallens/SocrateAI-Scientific-AlienMathematics-Foundation}

% ===================================================================
% ACKNOWLEDGMENTS
% ===================================================================
\subsection*{Acknowledgments}
The Lean~4 verification uses the Mathlib library. The author thanks
the Lean and Mathlib communities for their foundational work on
formalized mathematics.

% ===================================================================
% BIBLIOGRAPHY
% ===================================================================
\begin{thebibliography}{99}

\bibitem{AW2021}
J.~Alman and V.~Vassilevska~Williams.
\newblock A refined laser method and faster matrix multiplication.
\newblock In \emph{Proc.\ 32nd ACM-SIAM Symposium on Discrete
Algorithms (SODA)}, pages 522--539, 2021.

\bibitem{AlphaEvolve2025}
DeepMind AlphaEvolve Team.
\newblock Discovering faster matrix multiplication algorithms with
AlphaEvolve.
\newblock \emph{Nature}, 2025.
\newblock Achieved rank 48 for $\langle 4,4,4 \rangle$, verified
by symbolic computation.

\bibitem{Bini1979}
D.~Bini, M.~Capovani, F.~Romani, and G.~Lotti.
\newblock $O(n^{2.7799})$ complexity for $n \times n$ approximate
matrix multiplication.
\newblock \emph{Information Processing Letters}, 8(5):234--235, 1979.

\bibitem{BlaserLysikov2020}
M.~Bl\"aser and V.~Lysikov.
\newblock On the tensor rank of algebras of small rank.
\newblock \emph{Computational Complexity}, 29:article~7, 2020.

\bibitem{Blaser2003}
M.~Bl\"aser.
\newblock On the complexity of the multiplication of matrices of small
formats.
\newblock \emph{J.~Complexity}, 19(1):43--60, 2003.
\newblock Proved $R(\langle n,n,n \rangle) \ge \tfrac{5}{2}n^2 - 3n$.

\bibitem{Bini1980}
D.~Bini.
\newblock Relations between exact and approximate bilinear algorithms.
Applications.
\newblock \emph{Calcolo}, 17(1):87--97, 1980.
\newblock Connects border rank (degeneration) to approximate bilinear
complexity.

\bibitem{BrydgesSpencer1985}
D.~Brydges and T.~Spencer.
\newblock Self-avoiding walk in 5 or more dimensions.
\newblock \emph{Comm.\ Math.\ Phys.}, 97(1--2):125--148, 1985.

\bibitem{CampanelliSimmonsDuffin2024}
M.~Campanelli and D.~Simmons-Duffin.
\newblock The critical exponent $\gamma$ for self-avoiding walks:
a conformal bootstrap approach.
\newblock Preprint, 2024.

\bibitem{Clisby2010}
N.~Clisby.
\newblock Accurate estimate of the critical exponent $\nu$ for
self-avoiding walks via a fast implementation of the pivot algorithm.
\newblock \emph{Physical Review Letters}, 104(5):055702, 2010.

\bibitem{CW1990}
D.~Coppersmith and S.~Winograd.
\newblock Matrix multiplication via arithmetic progressions.
\newblock \emph{Journal of Symbolic Computation}, 9(3):251--280, 1990.

\bibitem{DWZ2023}
R.~Duan, H.~Wu, and R.~Zhou.
\newblock Faster matrix multiplication via asymmetric hashing.
\newblock In \emph{Proc.\ 64th Annual IEEE Symposium on Foundations
of Computer Science (FOCS)}, pages 2129--2138, 2023.

\bibitem{Hammersley1957}
J.~M.~Hammersley.
\newblock Percolation processes: II.\ The connective constant.
\newblock \emph{Proc.\ Cambridge Philos.\ Soc.}, 53(3):642--645, 1957.

\bibitem{HillarLim2013}
C.~J.~Hillar and L.-H.~Lim.
\newblock Most tensor problems are NP-hard.
\newblock \emph{J.~ACM}, 60(6):article~45, 2013.

\bibitem{HaraSlade1992}
T.~Hara and G.~Slade.
\newblock Self-avoiding walk in five or more dimensions.\ I.\ The
critical behaviour.
\newblock \emph{Comm.\ Math.\ Phys.}, 147(1):101--136, 1992.

\bibitem{LandsburgManivel2008}
J.~M.~Landsberg and L.~Manivel.
\newblock Generalizations of Strassen's equations for secant varieties
of Segre varieties.
\newblock \emph{Comm.\ Algebra}, 36(2):405--422, 2008.

\bibitem{Pan1980}
V.~Ya.~Pan.
\newblock New fast algorithms for matrix operations based on
Strassen's algorithm.
\newblock \emph{SIAM J.\ Comput.}, 9(2):321--342, 1980.

\bibitem{Schonhage1981}
A.~Sch\"onhage.
\newblock Partial and total matrix multiplication.
\newblock \emph{SIAM J.\ Comput.}, 10(3):434--455, 1981.

\bibitem{SocrateAI2026}
X.~Callens.
\newblock SocrateAI Scientific Agora: Verified Lean~4 modules for
non-commutative algebra and tensor decomposition.
\newblock \url{https://github.com/xaviercallens/SocrateAI-Scientific-AlienMathematics-Foundation},
2026.

\bibitem{Strassen1969}
V.~Strassen.
\newblock Gaussian elimination is not optimal.
\newblock \emph{Numer.\ Math.}, 13:354--356, 1969.

\end{thebibliography}

\end{document}
